\subsubsection*{3.1.9.}
参见 [3]。设 $X$ 是拓扑空间，$i:Y\hookrightarrow X$ 是闭子集，$T$ 是 $X$ 的自同构，在 $Y$ 上诱导恒等映射，并设 $\Lambda$ 是交换群。为定义变差 $\operatorname{Var}$，选取 $X$ 上 $\Lambda$ 的一个 $T$-等变消解 $\Lambda_X^*$，以及 $Y$ 上常值层 $\Lambda$ 的一个消解 $\Lambda_Y^*$，约定 $T$ 在其上平凡作用；再选取一个等变态射
\[
\varphi:\Lambda_X^*\longrightarrow i_*\Lambda_Y^*
\]
消解显然的态射 $\varphi_0:\Lambda_X\to i_*\Lambda_Y$。假设层 $\Lambda_X^i$ 与 $\Lambda_Y^i$ 对所考虑的函子均无上同调。取整体截面得
\[
\Gamma(\varphi):\Gamma(\Lambda_X^*)
\longrightarrow\Gamma(\Lambda_Y^*).
\]
为简单起见，假设每个 $\Gamma(\varphi)^i$ 均满；否则，可通过向 $\Gamma(\Lambda_X^*)$ 添加 $\Gamma(\Lambda_Y^*)$ 的锥使其为满。令
\[
R\Gamma(X\bmod Y,\Lambda)=\operatorname{Ker}(\Gamma(\varphi)).
\]
在合理的假设下，若 $\mathfrak F$ 是 $X$ 中与 $Y$ 不交的闭子集构成的族，则有拟同构
\[
\Gamma_{\mathfrak F}(X-Y,\Lambda_X^*)
\xrightarrow{\ \sim\ }R\Gamma(X\bmod Y,\Lambda).
\]
$\Gamma(\Lambda_X^*)$ 的自同态 $T-1$（其中 $T$ 通过搬运结构作用）在 $\Lambda$（常值部分）上为零，并经 $\operatorname{Ker}(\Gamma(\varphi))$ 分解。因此，在导出范畴中且在相同假设下，它给出变差态射
\[
\operatorname{Var}:R\widetilde\Gamma(X,\Lambda)
\longrightarrow R\Gamma_{\mathfrak F}(X-Y,\Lambda).
\]
无这些假设时，则得到
\[
\operatorname{Var}:R\widetilde\Gamma(X,\Lambda)
\longrightarrow R\Gamma(X\bmod Y,\Lambda).
\]

当 $T$ 在 $Y$ 的一个邻域上为恒等，并在该邻域内于 $\Lambda_X^*$ 上平凡作用时，此描述可简化。此时 $\Gamma(\Lambda_X^*)$ 的自同态 $T-1$ 经 $\Gamma_{\mathfrak F}(\Lambda_X^*)$ 分解，并给出变差。

\subsubsection*{命题 3.1.10.}
在 3.1.8(A),(B) 的同构下，变差
\[
(R\Phi_\eta(\Lambda))_0
\longrightarrow R\Gamma_{\{0\}}(R\Psi_\eta(\Lambda))
\]
等于 3.1.9 中由 $V$ 的自同构 $T$ 所定义的变差，此自同构在 $\partial V$ 上为恒等。

建议读者画一张小图来理解，而不必阅读证明。设 ${}_fB$ 是 3.1.2 的闭盒形区域，${}_fC\subset{}_fB$ 是沿其边界的领圈；于是通过将 $\partial_fB\times\{1\}$ 恒等映至 $\partial_fB$ 的一个同构，${}_fC\simeq\partial_fB\times[0,1]$ 成为 $D_t$ 上的丛。设
\[
{}_f\alpha:{}_fC\xrightarrow{\ \sim\ }
(\partial V\times[0,1])\times D_t
\]
为 $D_t$ 上丛 ${}_fC$ 的平凡化，其在 $y\in D_t$ 处的纤维诱导恒等映射
\[
\partial V=\partial(f^{-1}(y))
\xrightarrow{\ \sim\ }\partial V\times\{1\}.
\]
令 $C={}_fC\cap B$，并设
\[
\alpha:C\xrightarrow{\ \sim\ }
(\partial V\times[0,1))\times D_t
\]
由 ${}_f\alpha$ 诱导。保留 (3.1.8.2) 的记号。

设 $R$ 是 $\overline B$ 上由以下规则定义的等价关系：$xRy$，当且仅当 $x=y$，或者 $p(x),p(y)\in C$ 且 $\alpha(p(x))$、$\alpha(p(y))$ 在 $\partial V\times[0,1)$ 中有相同的像。令
\[
{}_1\overline B=\overline B/R,
\qquad
{}_1\overline B^*=\overline B^*/R,
\qquad
{}_1\overline C^*=p^{-1}(C)/R.
\]
有交换图
\[
\begin{tikzcd}[column sep=large,row sep=large]
B_0 \arrow[r,hook,"\bar\imath"] \arrow[d,equal]
  & \overline B \arrow[d,"q"]
  & \overline B^* \arrow[l,hook',"\bar\jmath"'] \arrow[d,"q'"] \\
B_0 \arrow[r,hook,"\bar\imath"]
  & {}_1\overline B
  & {}_1\overline B^* \arrow[l,hook',"{}_1\bar\jmath"']
\end{tikzcd}
\]
设 $\Lambda^*$ 是 $\overline B^*$ 上由无上同调层构成的 $\Lambda$ 的一个 $\pi_1(D^*)$-等变消解。按定义，变差由复形态射
\[
\Lambda[0]\longrightarrow
\bar\imath^*\bar\jmath_*(\Lambda^*).
\]
计算。由于 $\Lambda\xrightarrow{\sim}Rq_*\Lambda$ 且 $\Lambda\xrightarrow{\sim}Rq'_*\Lambda$，也可使用 $\Lambda$ 在 ${}_1\overline B^*$ 上的一个等变消解及态射
\[
\Lambda[0]\longrightarrow
\bar\imath^*{}_1\bar\jmath_*(\Lambda).
\]
来计算。在 (A) 和 (B) 下，变差由此等同于复合态射
\begin{equation}\tag{3.1.10.1}
\begin{tikzcd}[column sep=small,row sep=large,cells={nodes={font=\small}}]
R\Gamma(V,\Lambda) \arrow[r,"\sim"]
  & R\Gamma(V^\circ,\Lambda)
  & R\Gamma({}_1\overline B^*,\Lambda)
      \arrow[l,"\sim"',"\text{\textcircled{1}}"]
      \arrow[d,"\operatorname{Var}"] \\
{}
  & R\Gamma_c(V^\circ,\Lambda)
  & R\Gamma({}_1\overline B^*\bmod{}_1\overline C^*,\Lambda)
      \arrow[l,"\sim"',"\text{\textcircled{2}}"]
\end{tikzcd}
\end{equation}
其中 \textcircled{1} 与 \textcircled{2} 由 $V^\circ$ 到 ${}_1\overline B^*$ 的包含诱导，且 $\operatorname{Var}$ 是 3.1.9 对 ${}_1\overline B^*$ 及 $\pi_1(D^*)$ 的正生成元的作用所定义的变差。此群通过 1.1.5 中描述的作用来作用于 ${}_1\overline B^*$，并平凡地作用于 ${}_1\overline C^*$。

在 $B^*/D^*$ 上选择与 $C$ 的平凡化相容的联络 $\nabla$。此联络平凡化 $\overline B^*$，取商后定义映射 $u:{}_1\overline B^*\to V$。对如 3.1.3 中定义的 $T$，图
\[
\begin{tikzcd}[column sep=large,row sep=large]
{}_1\overline B^* \arrow[r,"u"] \arrow[d,"T"']
  & V \arrow[d,"T"] \\
{}_1\overline B^* \arrow[r,"u"]
  & V
\end{tikzcd}
\]
交换。映射 $u$ 诱导上同调同构，故命题 3.1.10 由 (3.1.10.1) 及变差的函子性得证。

命题 3.1.10 是为后面的一个显式计算提供依据；该计算中同此处一样，$T$ 在 $\partial V$ 的一个邻域上为恒等。我们将取一个由流 $C$ 表示的上同调类 $c$，并使用如下事实：$\operatorname{Var}(c)$ 由在 $V^\circ$ 中具有紧支集的流 $TC-C$ 表示。

